% Transcription — fonds Alexandre Grothendieck, shelfmark 21, batch 2 (pages 21-28).
% Established from the facsimile archives/batches/21/batch-02.pdf (a mirror of
% https://grothendieck.umontpellier.fr/21.pdf), per the skill
% transcribe-grothendieck. The batch file carries no cover sheet and opens at the
% archivists' page 21: PDF page N is archivists' page N + 20. Checked against the
% pencilled numbers bottom-left — "22", "23", "24", "25", "26", "27", "28" read on
% the corresponding leaves — which confirm the alignment end to end. The folder has
% 28 pages, so this batch is its close.
%
% Pass: Opus 5 (claude-opus-5), 2026-09-19 — first pass, unchecked against the pages by a human.
%
% What the batch is. Eight leaves, ALL OF THEM IN HIS HAND, in ink over a lighter
% ink he goes back over: this is the "notes manuscrites (s.d.)" half of the
% inventory's title, not the "tapuscrit et copies de tapuscrit annotés" half. There
% is no typed leaf in the batch, no overtyping, no interlinear annotation of a typed
% page — so the apparatus of an annotated typescript (\struck for xxxx'd-out words,
% \add for interlinear hand, \marginal for hand in the margin of a typed page) has
% almost no occasion here: the one \marginal is a note he writes up the left edge of
% his own manuscript leaf, and every \struck is his pen through his own ink.
%
% The run is a single continuous argument, unpaginated by him, carried by numbered
% statements: Définition 6 (p. 21), Prop. 7 and Corollaire 8 (p. 24), Corollaire 10
% (p. 27), Proposition 9 (p. 28). The numbering is his draft's and it is out of
% order at the end — 8, then 10, then 9 — and it is transcribed as written. It does
% not correspond to the numbering of any printed EGA and nothing here has been
% checked against one.
%
% The subject is the extension of objects and of homomorphisms across a closed
% subset — Lefschetz-type conditions, but stated for a FIBERED CATEGORY rather than
% for a sheaf:
%
%   21   the tail of a proof from p. 20, then Définition 6: X a topological space,
%        Y ⊂ X closed, 𝒰 = X - Y, 𝓕 a fibered category over the opens of X. 𝓕
%        satisfies P₁ (P₂, P₃) rel. to Y when 𝓕(V) → 𝓕(V ∩ 𝒰) is faithful (fully
%        faithful, an equivalence) for every open V; and at a point x when the same
%        holds after passing to the PSEUDO-LIMIT over the open neighbourhoods W of
%        x. A marginal note asks whether the global and pointwise forms agree on a
%        locally noetherian spectral space.
%   22   what the two forms say in terms of the sheaf Hom(f,g): for P₁ and P₂ the
%        conditions on 𝓕 are conditions on Hom(f,g) rel. to Y; P₃ adds essential
%        surjectivity — every object over V ∩ 𝒰 is the restriction of one over V.
%   23   a third form: introduce the "catégorie fibrée induite" 𝓕_(x) on X_(x) by
%        𝓕_(x)(W) = Pseudolim 𝓕(W); then Hom(f,g)_(x) ≃ Hom(f_(x), g_(x)), and the
%        N.B. observes that P₁ (P₂, P₃) for 𝓕 at x is P₁ (P₂, P₃) for 𝓕_(x) on X_(x).
%   24   Prop. 7 — for X a noetherian spectral space, 𝓕 satisfies P₁ (P₂, P₃) rel.
%        to Y iff it satisfies it at every x ∈ Y; the reduction of the proof to P₃
%        assuming P₁, P₂ already had; and Corollaire 8, the stability under passing
%        to an open X' ⊂ X and a closed Y' ⊂ Y.
%   25   the proof: a commutative triangle of restriction functors 𝓕(V) → 𝓕(V ∩ 𝒰'),
%        𝓕(V ∩ 𝒰'') → 𝓕(V ∩ 𝒰), two of whose legs are equivalences; then (i) ⟹ (ii)
%        by applying the condition to the closure x̄ of a point and passing to the
%        limit.
%   26   (ii) ⟹ (i) by noetherian induction on Y: take x maximal in Y, descend to
%        𝓕_(x), lift to some 𝓕(W), glue, and conclude by the induction hypothesis.
%   27   Corollaire 10, the translation into local cohomology: the conditions are
%        H^i_Y(F) = 0 in a range of i, the sheaf H_Y being described as the sheaf
%        associated to V ↦ H(V ∩ 𝒰, F).
%   28   the special case of TORSORS: F a sheaf of groups on X, 𝓕 the fibered
%        category of principal homogeneous spaces under F; Proposition 9 says that
%        P₁ (P₂) for 𝓕 rel. to Y, P₁ (P₂) for F rel. to Y, and the vanishing of
%        H⁰_Y(F) (and H¹_Y(F)) are equivalent — the proof identifying F with
%        Hom(f₀,f₀) for f₀ the trivial torsor, and reading H⁰_Y and H¹_Y as Ker and
%        Coker of F → i_* i^* F.
%
% Legibility. This is drafting, not copying out: the mathematical statements come
% through, the connective sentences between them very often do not, and the batch
% therefore carries a great many \ill{}. Nothing has been filled. Two recurring
% conventions of his are worth naming once: he writes "à" as a colon, so "rel. : Y"
% throughout is "rel. à Y" and is transcribed as he wrote it; and "tt" is "tout",
% which the guide warns reads as an H.
\documentclass[11pt,a4paper]{article}
% Le préambule commun à toutes les transcriptions du fonds Grothendieck.
%
% Il tient en une idée : **l'incertitude se compose, elle ne se résout pas.**
% Une transcription qui lisse un mot douteux a détruit la seule chose pour
% laquelle on la faisait — un lecteur doit pouvoir savoir, à chaque endroit,
% ce qui était sur le feuillet et ce qui a été ajouté. Les six macros qui
% suivent sont donc l'appareil critique, et non de la décoration.
%
% Les mêmes commandes sont reconnues par `scripts/render.mjs`, qui produit la
% vue de lecture du site. Une macro ajoutée ici doit l'être aussi là-bas,
% faute de quoi le rendu échouera bruyamment — ce qui est voulu : un rendu
% silencieusement faux est pire qu'un rendu absent.

\ProvidesPackage{grothendieck}[2026/08/08 Transcriptions du fonds Grothendieck]

\usepackage{amsmath,amssymb,amsthm}
\usepackage{mathrsfs}
\usepackage{stmaryrd}
% Les diagrammes commutatifs sont partout dans ce fonds : c'est la langue
% même de la géométrie algébrique de Grothendieck, et une transcription qui
% les rendrait en prose ne transcrirait rien.
\usepackage{tikz-cd}
% Français par défaut — c'est la langue des feuillets — mais l'anglais est
% chargé aussi : la traduction et le résumé sont des documents anglais, et la
% typographie française (espaces avant les deux-points) y serait une faute.
% Ces documents basculent par \selectlanguage{english} en tête de corps.
\usepackage[english,french]{babel}
\usepackage{fontspec}
\usepackage{geometry}
\geometry{a4paper,margin=25mm,marginparwidth=28mm}
\usepackage{marginnote}
\usepackage{xcolor}
% ulem, not soul. `soul`'s \st and \ul are built on a character-by-character
% scan that XeLaTeX's font machinery breaks: the document fails with an
% undefined control sequence rather than degrading. ulem does the same two jobs
% and is Unicode-engine safe. `normalem` stops it hijacking \emph, which the
% transcriptions use for Grothendieck's own emphasis.
\usepackage[normalem]{ulem}
\usepackage{hyperref}
\hypersetup{colorlinks=true,linkcolor=black,urlcolor=black,pdfborder={0 0 0}}

% --- Métadonnées du lot -----------------------------------------------------
% Reprises telles quelles par le script de rendu, qui les affiche en tête de
% la vue de lecture. Elles fixent ce que le fichier prétend couvrir : sans
% elles, une transcription flotte sans point d'attache dans un fonds de seize
% mille feuillets.

\newcommand{\folder}[1]{\gdef\grothFolder{#1}}
\newcommand{\batch}[1]{\gdef\grothBatch{#1}}
\newcommand{\leaves}[2]{\gdef\grothFirst{#1}\gdef\grothLast{#2}}
\newcommand{\foldertitle}[1]{\gdef\grothFoldertitle{#1}}
\gdef\grothFolder{?}\gdef\grothBatch{?}\gdef\grothFirst{?}\gdef\grothLast{?}\gdef\grothFoldertitle{}

% --- Le résumé --------------------------------------------------------------
%
% Il ouvre la lecture modernisée, sous le titre « Résumé », et il est composé
% en petit. Ce n'est pas un ornement : c'est le seul passage du document qui
% ne soit pas des mathématiques, et un lecteur doit pouvoir le reconnaître
% avant de l'avoir lu — puis le sauter s'il connaît déjà le sujet, ou s'y
% arrêter s'il ne le connaît pas. Le filet qui le suit marque la reprise du
% corps.

\newenvironment{resume}
  {\small\setlength{\parskip}{0.45em}}
  {\par\medskip\hrule\medskip}

% --- Mots-clés ---------------------------------------------------------------
%
% En anglais, en clôture du résumé de la lecture modernisée : le vocabulaire
% moderne sous lequel chercher ce que les feuillets construisent. Le manifeste
% du site (`npm run manifest`) les extrait pour étiqueter la cote entière —
% cette ligne est la source unique des tags, il n'y a pas de fichier à côté.
\newcommand{\keywords}[1]{\par\smallskip\noindent
  {\footnotesize\textsc{Keywords} --- \textit{#1}}\par}

% --- L'appareil critique ----------------------------------------------------

% Le numéro de feuillet, en marge. C'est l'ancre qui permet de régler un
% désaccord sur une formule en regardant *un* feuillet plutôt qu'une chemise
% de six cents. Le site s'en sert aussi pour tourner les pages du fac-similé
% quand on fait défiler la transcription.
\newcommand{\leaf}[1]{%
  \marginnote{\footnotesize\color{gray!70}\textsf{#1}}%
  \label{leaf:#1}\ignorespaces}

% Illisible : on ne devine pas. Un mot inventé qui se lit comme les autres est
% le pire résultat possible.
\newcommand{\ill}{\textcolor{red!60!black}{[\,\ldots\,]}}

% Lecture douteuse : le mot est donné, et signalé comme incertain.
\newcommand{\uncertain}[1]{\uline{#1}}

% Ajout de l'éditeur — une lettre manquante, un mot sous-entendu.
\newcommand{\add}[1]{\textcolor{blue!50!black}{[#1]}}

% Biffé par Grothendieck lui-même. Ce qu'il a rayé fait partie du document :
% une rature dit souvent où la pensée a changé de direction.
\newcommand{\struck}[1]{\sout{#1}}

% Note du transcripteur, sur l'état matériel du feuillet.
\newcommand{\note}[1]{\par\smallskip{\small\color{gray!60!black}\textit{[#1]}}\par\smallskip}

% Note marginale de Grothendieck, distincte du corps du texte.
\newcommand{\marginal}[1]{\par\smallskip\noindent\hspace{1em}%
  {\small\color{gray!60!black}#1}\par\smallskip}

% --- Titre ------------------------------------------------------------------

\AtBeginDocument{%
  \begin{center}
    {\large\bfseries Fonds Alexandre Grothendieck --- cote n\textsuperscript{o}~\grothFolder}\\[2pt]
    {\itshape\grothFoldertitle}\\[4pt]
    {\small feuillets \grothFirst--\grothLast\ (lot \grothBatch)}\\[2pt]
    {\footnotesize\color{gray!60!black}Transcription établie d'après le fac-similé numérisé
     par l'Université de Montpellier.\\
     Les lectures incertaines sont soulignées, les illisibles notées [\,\ldots\,],
     les ajouts entre crochets.}
  \end{center}
  \vspace{1em}}

\endinput


\folder{21}
\batch{2}
\pages{21}{28}
\dating{[vers 1960-1967]}
\watermark{Édition de démonstration}
\foldertitle{[Chapitre] 0 III : tapuscrit et copies de tapuscrit annotés (s.d.), notes manuscrites (s.d.).}

\begin{document}

\section*{Les conditions $P_1$, $P_2$, $P_3$ pour une catégorie fibrée}

\note{le titre est de l'éditeur : le feuillet 21 reprend en cours de
démonstration, sans en-tête, la suite du feuillet 20 (batch 1)}

\page{21}
$f \in \Gamma(\mathcal{U}, F)$, \ill{} on \uncertain{peut} prolonger $f$
en une section de $\mathcal{F}$ sur $X$. Il s'agit de prolonger \ill{}
un voisinage de $x$, \ill{} l'hyp.\ de récurrence. \ill{} c'est encore
possible \ill{} l'hyp.\ $P_2$ en $x$, et la \ill{}. OK.

\note{le début du feuillet achève une démonstration commencée au feuillet
précédent, qui n'est pas dans ce batch}

\emph{Définition 6.} Soit $\mathcal{F}$ une catégorie fibrée \ill{} sur la
catégorie des ouverts d'un \uncertain{espace} topologique $X$, soit $Y$ une
partie fermée de $X$. On dit que $\mathcal{F}$ satisfait la condition $P_1$
(resp.\ $P_2$, resp.\ $P_3$) rel.\ à $Y$, si pour tout ouvert $V \subset X$,
le foncteur
\[
  \Gamma\mathcal{F}(V) \longrightarrow \mathcal{F}(V \cap \mathcal{U})
  \qquad (\mathcal{U} = X - Y)
\]
est fidèle (resp.\ pleinement fidèle, resp.\ une équivalence de catégories).

\note{le terme de gauche porte devant le $\mathcal{F}$ un jambage
supplémentaire, lu $\Gamma$, que le terme de droite n'a pas ; l'asymétrie est
sur le feuillet et la lecture reste douteuse}

On dit que $\mathcal{F}$ satisfait la condition $P_1$ (resp.\ $P_2$, resp.\
$P_3$) \emph{en un point} $x \in X$, si le foncteur
\[
  \mathrm{Pseudolim}_{W}\ \mathcal{F}(W) \longrightarrow
  \mathrm{Pseudolim}_{W}\ \mathcal{F}(W - W \cap \bar{x})
\]
\ill{} est fidèle (resp.\ pleinement fidèle, resp.\ une équivalence).

\note{sur le feuillet, le $W$ de chaque \emph{Pseudolim} est écrit sous une
longue flèche placée au-dessous du mot ; $W$ parcourt les voisinages ouverts
de $x$}

\marginal{\ill{} loc.\ \uncertain{noethérien} \uncertain{spectral} \ill{},
$\mathcal{F}$ rel.\ \ill{} $\Longrightarrow$ la condition $P_1$ ($P_2$, $P_3$)
\ill{} ; \ill{} la condition $P_1$ ($P_2$, $P_3$) rel.\ \ill{} $\{x\}$,
\ill{} $x$ \ill{} $V$ \uncertain{fermé}.}

\page{22}
La première \ill{} signifie, pour $P_1$ (resp.\ $P_2$) que pour tt ouvert
$V \subset X$, et tt couple $f, g \in \mathrm{Ob}\ \mathcal{F}(V)$, le
faisceau $\underline{\mathrm{Hom}}(f,g)$ sur $V$ satisfait la condition $P_1$
(resp.\ $P_2$) rel.\ à \struck{\ill{}} $Y \cap V$. La condition $P_3$ signifie
que de plus, pour tt ouvert $V \subset X$, et tt objet
$f' \in \mathrm{Ob}\ \mathcal{F}(V \cap \mathcal{U})$ est isomorphe à un objet
de la forme $f | V \cap \mathcal{U}$ où $f \in \mathcal{F}(W)$.

\note{le dernier terme est écrit $\mathcal{F}(W)$ alors que l'argument demande
$\mathcal{F}(V)$ ; transcrit tel quel}

La seconde \ill{} signifie que le foncteur
\[
  \mathcal{F}(X) \longrightarrow \mathcal{F}(X - \{x\})
\]
\ill{} fidèle \struck{\ill{}} (resp.\ \struck{fidèle}, une équivalence)
\ill{} \struck{$P_1$ ($P_2$)} (pour tt $f, g \in \mathrm{Ob}\ \mathcal{F}(X)$,
le faisceau $\underline{\mathrm{Hom}}(f,g)$ satisfait $P_1$ ($P_2$), et pour
$P_3$, que tt objet $f'$ de $\mathcal{F}(X - \{x\})$ soit isom.\ à un objet de
la forme $f | X - \{x\}$, où $f \in \mathrm{Ob}\ \mathcal{F}(X)$.

\note{la parenthèse ouverte avant « pour tt $f,g$ » n'est jamais refermée sur
le feuillet}

\page{23}
La troisième \ill{} se \uncertain{ramène} \ill{} en général, si on
\uncertain{introduit} pour tt $x \in X$ la « catégorie fibrée induite »
$\mathcal{F}_{(x)}$, par la condition que pour tt \struck{\ill{}} $W$ de
$X_{x}$, \ill{} que
\[
  \mathcal{F}_{(x)}(W) \;=\; \mathrm{Pseudolim}_{W}\ \mathcal{F}(W)
\]
où $W$ parcourt les ouverts de $X$ tels que $W_{(x)} \supset \struck{W}\ w$
(ou \ill{} \uncertain{ouvert}, tel que $W_{(x)} = w$).
\struck{\ill{}}

\ill{} cette définition, pour tt \ill{} $V$ de $X$ et \ill{}
$f, g \in \mathrm{Ob}\ \mathcal{F}(V)$, \ill{} \struck{\ill{}} isom.\ (ou
\[
  \underline{\mathrm{Hom}}(f,g)_{(x)} \;\Longrightarrow\;
  \underline{\mathrm{Hom}}(f_{(x)}, g_{(x)})
\]
\ill{} des définitions.

\note{le feuillet porte ici une double flèche, non une flèche d'isomorphisme ;
transcrite telle quelle}

\emph{N.B.} La propriété $P_1$ ($P_2$, $P_3$) de $\mathcal{F}$ en $x$,
\uncertain{j'imagine} \ill{} revient à dire que $\mathcal{F}_{(x)}$ sur
$X_{(x)}$ satisfait à la propriété $P_1$ ($P_2$, $P_3$). \ill{} que la
propriété $P_1$ ($P_2$) est équivalente \ill{} que pour \struck{tt ouvert
$V$} \struck{de $x$}, et $f, g \in \mathrm{Ob}\ \mathcal{F}(W)$,
$\underline{\mathrm{Hom}}(f,g)$ satisfait à $P_1$ ($P_2$) en $x$.

\page{24}
On voit donc, grâce à \ill{} la validité des \ill{} \struck{\ill{}}
\struck{dans}

\emph{Prop.\ 7.} Soit $X$ un esp.\ top.\ \uncertain{noeth.} \uncertain{spectral},
$\mathcal{F}$ une catégorie fibrée \ill{} sur la catégorie des ouverts de $X$,
et soit $Y$ une partie fermée de $X$. \struck{\ill{}} Conditions
équivalentes :
\begin{enumerate}
  \item[(i)] $\mathcal{F}$ satisfait à $P_1$ ($P_2$, $P_3$) rel.\ à $Y$ ;
  \item[(ii)] pour tt $x \in Y$, $\mathcal{F}$ satisfait à $P_1$ ($P_2$,
        $P_3$) en $X$.
\end{enumerate}

\note{(ii) porte un $X$ majuscule là où l'énoncé demande le point $x$ ;
transcrit tel quel}

Reste le cas $P_3$, \struck{\ill{}} quand $P_1$, $P_2$ sont déjà vérifiés.
\ill{} On va prouver le
\[
  \text{(ii)} \Longrightarrow \text{(i)}
\]

\emph{Corollaire 8.} Si $\mathcal{F}$ satisfait $P_1$ ($P_2$, $P_3$) rel.\ à
$Y$, alors \struck{tt \ill{} $X' \subset X$ et} tt \ill{} $Y' \subset$
\struck{$Y$} \ill{} toute partie \ill{} \struck{$Y' \subset Y \cap X'$},
\struck{$\mathcal{F} | X'$} satisfait $P_1$ ($P_2$, $P_3$) rel.\ à $Y'$.

\struck{Déjà} fait pour $P_1$, $P_2$ ; voyons $P_3$. Voyons la \ill{}
\struck{\ill{}}

\page{25}

\begin{tikzcd}
  \mathcal{F}(V) \arrow[r, "\rho'"] \arrow[dr, "\rho"'] &
  \mathcal{F}(V \cap \mathcal{U}') \arrow[d, "\rho''"] \\
  & \mathcal{F}(V \cap \mathcal{U})
\end{tikzcd}

\ill{} $\rho''$ et $\rho$ sont des équivalences de catégories, donc $\rho'$ en
est une, OK.

Prouvons (i) $\Longrightarrow$ (ii), en appliquant la \uncertain{cond.} à
$\bar{x}$ (pour $x \in Y$). Soit $W$ un voisinage ouvert de $\bar{x}$, \ill{}
\[
  \mathcal{F}(W) \longrightarrow \mathcal{F}(W - W \cap \bar{x})
\]
une équivalence, donc \struck{a fortiori il} \ill{} pour $\mathcal{H}$ l'un
des types :
\[
  \mathcal{H}(\mathcal{F}(W)) \;\xrightarrow{\ \sim\ }\;
  \mathcal{H}(\mathcal{F}(W - W \cap \bar{x}))
\]
et \ill{} à la limite \ill{} :
\[
  \mathcal{H}(\mathcal{F}_{(x)}(X_{x})) \;\xrightarrow{\ \sim\ }\;
  \mathcal{H}(\mathcal{F}_{x}(X'_{\alpha_1}))
\]
O.K.

\note{la lettre $\mathcal{H}$ n'est pas introduite sur le feuillet ; elle y
tient lieu de « l'un des types » d'énoncé — fidélité, pleine fidélité,
équivalence}

Prouvons (ii) $\Longrightarrow$ (i). \ill{} pour \ill{} que tt
$f' \in \mathrm{Ob}\ \mathcal{F}(\mathcal{U})$ provient

\page{26}
d'un $f \in \mathrm{Ob}\ \mathcal{F}(X)$. À cause \ill{} déjà acquis, la
question est locale sur $X$, \ill{} $X$ \uncertain{noethérien}. Par suite, on
peut \ill{} (par récurrence \uncertain{noethérienne}) que pour tt
$Y' \subset Y$ et $\ne Y$, \struck{\ill{}} $\mathcal{F}$ satisfait $P_3$
rel.\ à $Y'$.

\struck{Soit} Si $Y = \emptyset$, il n'y a \ill{} rien à prouver ; \ill{},
soit $x$ un pt maximal de $Y$. Considérons $f' | X'_{(x)}$, par hypothèse il
provient d'un $f^{\circ}_{(x)} \in \mathrm{Ob}\ \mathcal{F}_{(x)}(X'_{(x)})$,
donc \ill{} provient d'un $f_0 \in \mathrm{Ob}\ \mathcal{F}(W)$, où $W$ est un
voisinage \ill{} de $x$. \struck{\ill{}} Considérons les restrictions de $f_0$
et $f'$ à $W \cap \mathcal{U} = V$, par construction elles deviennent
isomorphes sur $V_{(x)} = X'_{(x)}$ (donc \struck{\ill{}} par définition de
$\mathcal{F}_{(x)}$) à condition de prendre $W$ assez petit ; on peut supposer
qu'elles sont isomorphes. \ill{} les isomorphismes, on \ill{} un objet
$f'_1 \in \mathrm{Ob}\ \mathcal{F}(W \cup \mathcal{U})$ prolongeant $f'$. Par
l'hyp.\ de récurrence, $f'_1$ provient d'un $f \in \mathcal{F}(X)$, \ill{}

\page{27}
\ill{}

\emph{Corollaire 10.} Conditions équivalentes :
\begin{enumerate}
  \item[(i)] $F$ satisfait à $P_1$ ($P_2$, $P_3$) ;
  \item[(ii)] on a $\underline{H}^{i}_{Y}(F) = 0$ pour $i \leqslant 1$
        ($i \leqslant 2$, $i \leqslant 3$).
\end{enumerate}

Reste le cas $P_3$, \ill{}
\[
  \underline{H}^{2}_{Y}(F) \;=\; \text{faisceau associé au préfaisceau}
\]
\[
  V \rightsquigarrow H^{\ill}(V \cap \mathcal{U}, F).
\]

\note{l'exposant de ce dernier $H$ ne se lit pas — un seul jambage, ni un
zéro net ni un chiffre ; celui de $\underline{H}_Y$ se lit $2$}

On a (i) $\Longrightarrow$ (ii) puisque
$H^{\ill}(V \cap \mathcal{U}, F) \simeq H^{\ill}(V, F)$ (pour \ill{} $V$) et
que le faisceau associé \ill{} $V \rightsquigarrow H^{\ill}(V, F)$ \ill{}
tend \ill{}.

On a (ii) $\Longrightarrow$ (i). Car dans ce cas \ill{}
$f' \in \mathrm{Ob}\ \mathcal{F}(\mathcal{U})$ provient d'un
$f \in \mathrm{Ob}\ \mathcal{F}(X)$.
Or, \ill{} suffit de \ill{} \uncertain{constater}, i.e.\ le \ill{} sur $Y$,
\ill{} et \ill{} par Prop.\ 7, en \ill{} \ill{} un point quelconque
$x \in Y$.

\page{28}
\emph{Cas particulier.} Soit \struck{\ill{}} $F$ un préfaisceau \ill{}
\uncertain{de groupes} (ou un \ill{}) sur $X$. \ill{} la catégorie
$\mathcal{F}$ \ill{} possédant (\ill{}) \ill{} des fibrés principaux
homogènes \ill{} sous $F$. Soit $Y$ une partie fermée de $X$. Alors :

\note{la ligne où $\mathcal{F}$ est définie est très abrégée ; seule la
locution « fibrés principaux homogènes » et le rôle de $F$ y sont sûrs}

\emph{Proposition 9.} \struck{Les} Conditions équivalentes :
\begin{enumerate}
  \item[(i)] $\mathcal{F}$ satisfait $P_1$ ($P_2$) rel.\ à $Y$ ;
  \item[(ii)] $F$ satisfait $P_1$ ($P_2$) rel.\ à $Y$ ;
  \item[(iii)] on a $\underline{H}^{0}_{Y}(F) = 0$
        (\,$\underline{H}^{0}_{Y}(F) = 0$ et
        $\underline{H}^{1}_{Y}(F) = 0$\,)
        \struck{pour $i \leqslant 1$ ($i \leqslant 2$)} i.e.\ $i = 0$
        ($i = 0, 1$).
\end{enumerate}

En effet, \struck{on aura} si $f_0$ \ill{} le fibré \ill{} \ill{} sur $X$,
\ill{} $F \simeq \underline{\mathrm{Hom}}(f_0, f_0)$, d'où \ill{}
(ii) $\Longrightarrow$ (i). D'autre part, \ill{}
$f, g \in \mathrm{Ob}\ \mathcal{F}(V)$, $\underline{\mathrm{Hom}}(f,g)$ est
localement isomorphe (comme \ill{}) à $F$, d'où (i) $\Longrightarrow$ (ii).
Enfin (ii) $\Longleftrightarrow$ (iii), par définition de
$\underline{H}^{0}_{Y}(F)$ et $\underline{H}^{1}_{Y}(F)$ comme Ker et Coker de
\[
  F \longrightarrow i_{*} i^{*} F .
\]

\end{document}
